\section{Related Work}
\label{sec:related}

\paragraph{Passage-based retrieval-augmented generation.}
Retrieval-augmented generation conditions a language model on retrieved text instead of relying only on its parameters \citep{lewis2020rag}. In the standard setting the retrieved passages are concatenated in rank order, and the model is left to identify which sentences matter and how they combine. Two lines of work show that this presentation is not neutral. Models use information in the middle of a long context less reliably than information at its edges \citep{liu2024lost}, and the composition of the retrieved set, including irrelevant passages, changes generation quality in ways that retrieval scores do not predict \citep{cuconasu2024noise}. Both findings motivate our concern with how evidence is presented once it has been retrieved.

\paragraph{Finer retrieval units and compression.}
A second line of work changes the unit of evidence. Dense X Retrieval indexes propositions, i.e.\ atomic, self-contained factoids, and finds that proposition-level units improve retrieval and downstream QA under a fixed computation budget \citep{chen2024densex}; FActScore uses a similar atomic decomposition for evaluation \citep{min2023factscore}. Compression methods reduce the context instead: LongLLMLingua prunes prompts conditioned on the question \citep{jiang2024longllmlingua}, and ECoRAG compresses retrieved documents according to their evidentiality and retrieves more when the compressed text is insufficient \citep{jeong2025ecorag}. These methods treat the loss of information as an explicit trade-off against length. Our Protocol~B takes the same view, but holds the selection policy fixed across representations, so that the cost of a representation's formatting can be separated from the choice of content.

\paragraph{Graph-based and hierarchical retrieval.}
GraphRAG builds an entity graph with community summaries for query-focused summarisation \citep{edge2024graphrag}, RAPTOR organises chunks into a tree of recursive summaries \citep{sarthi2024raptor}, and HippoRAG and HippoRAG~2 use a knowledge graph with personalised PageRank as a long-term memory for multi-hop retrieval \citep{gutierrez2024hipporag,gutierrez2025hipporag2}. KG$^2$RAG links chunks through fact-level relations and organises the expanded set before generation \citep{zhu2025kg2rag}, LightRAG combines graph indexing with dual-level retrieval \citep{guo2024lightrag}, PathRAG retrieves and prunes relational paths and serialises them into the prompt \citep{chen2026pathrag}, and LinearRAG avoids relation extraction altogether, arguing that extracted relations are unreliable and costly, and builds a relation-free graph from entities and semantic links \citep{zhuang2026linearrag}. In these systems graph construction, retrieval and prompt formatting change together, so a difference in end accuracy cannot be attributed to the explicit relations shown to the model.

\paragraph{Analyses of when graphs help.}
Several recent studies ask when graph structure is worth its cost. Motivated by reports that GraphRAG frequently underperforms vanilla RAG on real-world tasks, GraphRAG-Bench evaluates the whole pipeline, from graph construction through retrieval to generation, on tasks ranging from fact retrieval to complex reasoning, and identifies the scenarios in which graph structures provide measurable benefits \citep{xiang2026whentographs}. A benchmark comparison of RAG and GraphRAG on question answering and query-based summarisation, under a unified evaluation protocol, finds that each has distinct strengths depending on the task and that combining them helps \citep{han2025ragvsgraphrag}, and a unified framework organises graph-based RAG methods so that they can be compared under the same experimental settings \citep{zhou2025indepth}. These studies compare \emph{systems}: when a graph method wins or loses, retrieval, indexing, the content that reaches the model and its serialisation all differ at once. Our study is narrower and complementary. We fix the source evidence to the gold passages, verify per question that every representation shows the same facts, and vary only the representation; under a budget we vary only how much of a fixed ranking each representation can afford. This isolates the effect of presentation and of serialisation overhead, which system-level comparisons cannot separate, at the price of not measuring retrieval benefits at all.

\paragraph{Multi-hop QA, verification and attribution.}
We evaluate on MuSiQue \citep{trivedi2022musique} and 2WikiMultiHopQA \citep{ho2020twowiki}, whose annotations provide supporting paragraphs and a decomposition into reasoning steps, and on HotpotQA \citep{yang2018hotpotqa}, HoVer \citep{jiang2020hover} and FEVER \citep{thorne2018fever} as secondary datasets. ALCE evaluates whether generated statements are supported by cited passages \citep{gao2023alce}; we record citations in the same spirit but do not interpret them as faithfulness measurements (\S\ref{sec:limitations}).

\paragraph{Relation to our earlier architecture.}
Our Experiment~0 includes a reconstruction of EG-RAG, an evidence-graph pipeline that decomposes passages into claims, predicts typed relations and selects a token-budgeted subgraph. We use it only as one retrieved-evidence baseline. Its components that the original description leaves unspecified are our own choices (Appendix~\ref{app:reconstruction}), and no result of the original architecture paper is reused here.
